% Isolated primary-source audit for the P09.2 product--Gauss packet.
% Not for live input before independent and production-Rethlas verification.

\paragraph{Final-scope addendum (26 August 2026).}
The frozen production statement replaces the provisional smoothness
hypothesis on (X) below by the natural hypothesis that (X) is an
algebraic (K)-variety.  Gubler's Proposition~7.2 is stated for an arbitrary
algebraic variety, and the final proof uses no smoothness of (X).  It also
defines the relative Gauss point for every (w\in N_{\mathbf R}), while the
constant-(K)-section density argument is invoked only for
(w\in N\otimes_{\mathbf Z}\Gamma).  Final independent source and proof
audits passed these exact quantifiers.  The provisional scope paragraph
below is retained as superseded audit history.

\paragraph{Scope and verdict.}
Let \(K\) be algebraically closed, complete, and nontrivially valued.
Let \(X\) be a smooth algebraic \(K\)-variety with an algebraic action of a
split torus \(T\), and let
\[
 h:X^{\mathrm{an}}\longrightarrow N_{\mathbf R},
 \qquad
 h(t\mathbin{\cdot}x)=\operatorname{trop}_T(t)+h(x),
\]
be integral PL.  For
\[
 F(w)=\tfrac12 b(w,w)+\ell(w)+c
\]
with \(b\) symmetric and nondegenerate, the primary sources support the
product--Gauss proof that \(F\circ h\) is an ordinary GK delta-weight if and
only if every \(\langle m,h\rangle\), \(m\in M\), is CLD-smooth, subject to
the scope firewalls below.  This audit does not itself verify the new
assembled theorem.

\paragraph{GK ambient condition.}
\cite[Definition~9.9(ii)]{GK17} requires a bidegree-\((1,1)\) superform on
the full ambient vector space \(N_{U,\mathbf R}\), with piecewise-smooth
coefficients, restricting to the facewise Hessian.
\cite[Definition~1.6, Remark~1.7, Example~2.10, and
Remark~9.10]{GK17} show that this is a codimension-zero delta-preform whose
coefficient functions are single-valued and continuous across walls.
Therefore the contraction argument may use one ambient coefficient tensor
\(G_0\) at the wall.  Arbitrary independent smooth Hessians on incident
cells would not suffice.  Algebraic pullback by
\(\mu:T\times X\to X\) is covered by \cite[Proposition~9.15]{GK17}.

\paragraph{Algebraization of analytic moments.}
\cite[Proposition~7.2]{Gub16} assumes the algebraically closed, complete,
nontrivially valued setting and, in its analytic section, strict analytic
spaces.  Given an analytic moment
\(\varphi:W\to T^{\mathrm{an}}\) on an open subset of an algebraic
analytification, it supplies, near each point, an algebraic moment
\(\varphi':U\to T\) on a very affine open such that
\[
                 \varphi_{\mathrm{trop}}
                 =\varphi'_{\mathrm{trop}}.
\]
It does not assert \(\varphi=\varphi'\).  A finite family is handled by
putting all coordinates in one product torus and applying the proposition
once.  Gubler's Section~4.12 then factors the resulting algebraic moment
through the canonical moment of the very affine open.  Equality of
tropicalizations is exactly what is needed to preserve weights and pull
ambient forms back by an integral-affine map.

\paragraph{Analytic local isomorphisms.}
GK Proposition~9.15 does not cover an arbitrary analytic local
isomorphism.  In the algebraizable scope the required invariance follows as
follows.  Pull the finitely many moment coordinates of a delta-presentation
back through the analytic isomorphism and adjoin the analytic unit relating
the pulled frame to an algebraic frame on the source.  Apply
\cite[Proposition~7.2]{Gub16} once to this product moment.  The new weight is
the old weight plus an affine coordinate, and the new ambient form is the
pullback of the old one.  Repeat for the inverse.  This proves invariance
between analytic opens of algebraic varieties for algebraic line bundles;
it makes no assertion for arbitrary nonalgebraizable analytic spaces.

\paragraph{Product Gauss point and the order of quantifiers.}
\cite[\S\S12.1.9--12.1.11]{CLD25} defines
\[
 z_w=\operatorname{sq}_x(w)\in T_{\mathscr H(x)}^{\mathrm{an}}
        \subset T\times X
\]
and proves continuity of \((w,x)\mapsto z_w\).
Hence one first fixes a delta-chart at \(z_{w_0}\), then fixes all moment
coordinates and all coefficient complexes, and only then moves \(w_0\)
slightly.  Every fixed chart unit is a quotient of finite Laurent
polynomials over \(\mathscr H(x)\).  Its nonunique-leading-monomial locus is
a finite union of proper affine hyperplanes.  Pullbacks of the finitely many
polyhedral walls give either proper affine walls or walls containing the
whole torus germ.  Since
\(\Gamma=-\log|K^\times|\) is nonzero, divisible, and dense, one may choose
\[
                 w\in N\otimes_{\mathbf Z}\Gamma
\]
outside all proper walls while remaining in the fixed chart.  Adjoining a
basis of torus characters makes the fibre-slope map
\(a:N_{\mathbf R}\hookrightarrow\mathbf R^s\) injective.  The minimal cell
\(C\) at the chosen point then satisfies
\[
                 a(N_{\mathbf R})\subset L_C,
\]
and the same holds on its star.  This is the precise noncircular
lineality argument.

\paragraph{Rational points near a rational Gauss point.}
Let \(H\) be complete with infinite residue field, let
\(T\simeq\mathbf G_{m,H}^r\), and let
\(w\in N\otimes\Gamma_H\).  Every analytic open neighbourhood of
\(\eta_w\) contains \(t_0\in T(H)\) with
\(\operatorname{trop}(t_0)=w\).
Choose \(c_i\in H^\times\) with \(-\log|c_i|=w_i\) and reduce to the unit
Gauss point.  A basic neighbourhood is controlled by finitely many analytic
functions \(f_j\) and intervals containing their Gauss norms.  Normalize
each \(f_j\); its initial reduction \(\widetilde f_j\) is a nonzero Laurent
polynomial.  Choose
\[
 \bar u\in(\widetilde H^\times)^r
       \setminus\bigcup_jV(\widetilde f_j),
\]
which is possible because \(\widetilde H\) is infinite, and lift its
coordinates to \(H^{\circ\times}\).  All \(f_j\) then attain their Gauss
norms at the lifted point.  Scaling by the \(c_i\) proves the claim.
For \(H=\mathscr H(x)\), the residue field contains the algebraically closed
field \(\widetilde K\), so it is infinite.  The condition
\(w\in N\otimes\Gamma_H\) is necessary: an irrational polyradius cannot be
the tropical value of an \(H\)-rational point.

\paragraph{Lifting the fibre point to a local section.}
By \cite[\S11.1.10]{CLD25}, an algebraic analytification is good, hence
almost good, and
\[
 \kappa(x)=\mathcal O_{X,x}/\mathfrak m_x
       \quad\text{is dense in}\quad\mathscr H(x).
\]
If an open \(\Omega\subset T\times X\) contains an
\(\mathscr H(x)\)-rational fibre point \((t_0,x)\), approximate the finitely
many nonzero coordinates of \(t_0\) by elements of \(\kappa(x)^\times\),
close enough that the resulting fibre point remains in \(\Omega_x\).
Lift them to germs \(g_i\in\mathcal O_{X,x}\), imposing
\[
                   |g_i(x)-t_{0,i}|<|t_{0,i}|.
\]
The germs are units.  On one common neighbourhood they define
\[
 s(y)=(g_1(y),\ldots,g_r(y),y).
\]
Since \(s(x)\in\Omega\), continuity and shrinking give \(s(U)\subset\Omega\).
This is also the approximation mechanism of
\cite[Lemma~12.2.3]{CLD25}.  Thus if
\[
 \langle m,\operatorname{trop}_T+h\rangle
\]
is smooth near \(z_w\), pullback by \(s\) and subtraction of
\(\sum_i m_i(-\log|g_i|)\) prove that \(\langle m,h\rangle\) is smooth near
\(x\).  No residue-line-bundle identification at \(z_w\) is used.

\paragraph{Raynaud sources.}
\cite[Proposition~19.2.5]{CLD25} gives, under
\(p^*L\simeq q^*M\), the exact metric-weight identity with a quadratic
function \(F_L\circ\tau\); its homogeneous degree-two part is
\(\frac12q_L\).  If \(q_L\) denotes the polarized bilinear form, the same
term is written \(\frac12q_L(w,w)\).
\cite[Theorem~B.6.4]{CLD25} says that ampleness of \(L\) forces \(q_L\)
positive definite, hence \(q_L^\sharp\) is an isomorphism.
\cite[Remark~19.2.3]{CLD25} says that all components of the canonical
tropicalization \(\tau\) are smooth if and only if the residual extension
of \(B_s\) by \(T_{\widetilde K}\) is isotrivial.  This is the terminal
criterion; no claim that a residue at one Gauss germ literally equals
\(Q_s\) is needed.

\cite[Remark~B.6.6]{CLD25} supplies algebraizability of \(E\) in the
abelian-variety case.  In the split case, \cite[\S19.2.1 and
\S\S B.5.12--B.5.13]{CLD25} describes \(p:E\to A^{\mathrm{an}}\) as the
analytic covering by the lattice.  The split hypothesis is essential:
\cite[\S B.5.14]{CLD25} warns that before splitting, \(p\) need not be a
local isomorphism.  The present algebraically closed setup is split.

\paragraph{Final firewall.}
The source chain passes this audit provided the proof:
\begin{enumerate}
\item fixes the chart and every wall before choosing the generic rational
 \(w\);
\item uses \(w\in N\otimes\Gamma\);
\item writes \(\kappa(x)\), rather than an unspecified local-ring field, as
 the dense subfield of \(\mathscr H(x)\);
\item derives analytic-local-isomorphism invariance from the simultaneous
 moment algebraization instead of attributing it to GK Proposition~9.15;
\item freezes whether \(q_L\) denotes a quadratic form or its polarized
 bilinear form; and
\item invokes CLD Remark~19.2.3 globally, without a one-point residue
 identification.
\end{enumerate}
The assembled result still requires hostile mathematical review and
production Rethlas verification before promotion.
